%%%%%%%%%%%%%%%%%%%%%%%%%%%%%%%%%%%%%%%%%%%%%%%%%%%%%%%%%%%%%%%%%%%%%%%%%%%%%%%
% WissKI User Meeting — book of abstracts (extended abstract) template
%
% Overleaf (import from GitHub):
%   New project → Import project → From GitHub → paste the repository URL
%   Main document: main.tex (Project menu → Main document, if needed)
%   Compiler: pdfLaTeX (Menu → Compiler)
%
% Local build: latexmk -pdf main.tex   or   make
%%%%%%%%%%%%%%%%%%%%%%%%%%%%%%%%%%%%%%%%%%%%%%%%%%%%%%%%%%%%%%%%%%%%%%%%%%%%%%%

\documentclass[a4paper,11pt]{article}

\usepackage[margin=2.5cm]{geometry}

% Prefer Noto Sans when available (Overleaf TeX Live); otherwise CM sans
\IfFileExists{noto.sty}{%
  \usepackage[sfdefault]{noto}%
}{%
  \renewcommand{\familydefault}{\sfdefault}%
}
\usepackage[T1]{fontenc}

\usepackage{graphicx}
\usepackage{booktabs}
\usepackage{url}
\usepackage{natbib}
\usepackage{hyperref}
\hypersetup{colorlinks=true, linkcolor=black, urlcolor=blue, citecolor=black}

\graphicspath{{img/}}

\title{Title of the Abstract for WissKI User Meeting}

\author{%
    First Author\textsuperscript{1}\thanks{Corresponding author: \url{email@example.com}}
    \and
    Second Author\textsuperscript{2}%
}

\date{%
    \vspace{-1em}
    \small\itshape
    \textsuperscript{1}Institution A, City, Country \\
    \textsuperscript{2}Institution B, City, Country%
}

\begin{document}

\maketitle

This is the extended abstract for the WissKI User Meeting. The entire document serves as the abstract, summarizing the main contributions and findings of the presented work.

\section{Introduction}
This is an introduction section demonstrating headings. WissKI (Wissenschaftliche KommunikationsInfrastruktur) is a great tool, and this template serves as a starting point for conference abstracts. You can include footnotes like this\footnote{This is a sample footnote.}.

\section{Main Content}
Here we demonstrate various elements requested for the template.

\subsection{Citations}
You can cite an article \cite{Abril07}, a book \cite{Kosiur01}, an internet document \cite{Thornburg01}, or a collection \cite{Spector90}.

\subsection{Block quotations}
For a passage of more than a few words, set it off with the \texttt{quote} environment (narrower margins, often used with slightly smaller type). Place the citation after the closing quotation marks, before the final period if your style requires it.

\begin{quote}\small
``Understanding policy-based networking requires a clear model of how rules are applied across distributed systems.'' \cite{Kosiur01}
\end{quote}

For a longer extract with one or more paragraphs, use \texttt{quotation} instead (first-line indentation for each paragraph):

\begin{quotation}\small
This is the first paragraph of a block quotation. It can run over several sentences.

This is a second paragraph in the same block. End the block with the source \cite{Spector90}.
\end{quotation}

\subsection{Inline quotations}
Short phrases in running text use \LaTeX{} ``curly'' opening and closing quotes (two backticks, two apostrophes). Example: Kosiur writes that ``policy-based networking'' ties configuration to business rules \cite{Kosiur01}.

For a quotation inside a quotation, switch to single quotes on the inner span: ``He replied, `not yet', and closed the issue.''

\subsection{Images and Figures}
Place image files under \texttt{img/} (see \texttt{README.md}). Uncomment the line below to include a real figure.

\begin{figure}[ht]
    \centering
    % \includegraphics[width=0.5\linewidth]{your-image}
    \rule{4cm}{3cm}
    \caption{A sample figure placeholder.}
    \label{fig:sample}
\end{figure}

\subsection{Tables}
Tables can be formatted with the \texttt{booktabs} package.

\begin{table}[ht]
    \centering
    \caption{A sample table showing data.}
    \label{tab:sample}
    \begin{tabular}{llr}
        \toprule
        \textbf{Component} & \textbf{Description} & \textbf{Quantity} \\
        \midrule
        WissKI & Core System & 1 \\
        Drupal & CMS Backend & 1 \\
        Ontology & Data Model & \textit{n} \\
        \bottomrule
    \end{tabular}
\end{table}

\subsection{Links and URLs}
With \texttt{hyperref} loaded, you can show a readable label while linking to a target: \href{https://wiss-ki.eu/}{WissKI project website}. For email, use \href{mailto:contact@example.com}{contact@example.com}.

To print the address itself (monospace, line-break friendly), use \verb|\url|: \url{https://wiss-ki.eu/documentation}. For a URL that should not become a clickable link in print (rare), you can use \texttt{\textbackslash nolinkurl\{...\}} from the \texttt{url} package: \nolinkurl{https://example.org/path}.

\section{Conclusion}
This concludes the sample template.

\bibliographystyle{plain}
\bibliography{references}

\end{document}
